\section{Geometric conventions and the two-point shift}\label{sec:geometry}
First apply an invertible shear $(x,y)\mapsto(x+ty,y)$ separating all
abscissae; only finitely many parameters $t$ fail. Convexity and general
position are preserved. Write $p<q$ when $x(p)<x(q)$, and
$m(p,q)=(y(q)-y(p))/(x(q)-x(p))$. A $k$-cup is an increasing sequence of
$k$ points whose consecutive slopes strictly increase; a cap has strictly
decreasing slopes. Length always counts vertices.

Every subsequence of a cup is a cup, and likewise for caps: a chord slope
is a positive weighted average of the slopes of the edges it spans.
A cup and cap sharing both endpoints have disjoint interiors, on opposite
sides of the common chord. Their union is in convex position and has the
sum of their lengths minus two vertices. These statements are used in
the planar category.

\begin{proposition}[Projective reduction]\label{prop:shift}
Let $Q_n$ be the maximum size of a set avoiding an $(n-1)$-cap and a convex $n$-gon. Then $\ES(n)\le Q_n+2$.
\end{proposition}
\begin{proof}
Take an $n$-gon-free set $S$ and a hull vertex $o$. After a translation
and rotation, $o=(0,0)$ and every other point has positive second
coordinate. For sufficiently small $t>0$, apply
\[
 T_t(x,y)=\left(\frac{x-t}{y+t^2},\frac1{y+t^2}\right).
\]
This projective map is invertible and has positive denominator on
$\conv(S)$, so it preserves convex position. The nonroot images stay
bounded, with distinct limiting first coordinates $x/y$: an equality
would make two original points collinear with $o$. The root becomes
$(-1/t,1/t^2)$. For nonroot images $(X_i,Y_i)<(X_j,Y_j)$,
\[
 \operatorname{orient}(T_t(o),T_t(p_i),T_t(p_j))
 =\frac{X_j-X_i}{t^2}+\frac{Y_j-Y_i}{t}+X_iY_j-Y_iX_j>0
\]
for small enough $t$. Thus the root is leftmost and all its triples
are cups. Deleting it leaves no $(n-1)$-cup; reflection in a horizontal
line gives a set counted by $Q_n$. Hence $|S|\le Q_n+1$ and
$\ES(n)\le Q_n+2$.
\end{proof}

\section{A common residual partition}\label{sec:partition}
We use the three-chain argument of Norin and Yuditsky~\cite[Claim~3.3]{NY}.

\begin{lemma}[Three-chain bridge]\label{lem:threechain}
Let $n\ge6$. Suppose an $(n-2)$-cup $F$ ends at $s$, an $(n-1)$-cup
$P$ starts at $s$ and ends at $t$, and an $(n-2)$-cap $J$ starts with
$(s,q)$. If $q\ne t$, their union contains an $(n-1)$-cap or a convex
$n$-gon. Other overlaps between these chains are permitted.
\end{lemma}
\begin{proof}
Let $f$ be the penultimate point of $F$, and let $v,u,t$ be the final
three points of $P$. If $(f,s,q)$ is a cap, prepend $f$ to $J$.
Otherwise $F,q$ is an $(n-1)$-cup and $m(f,s)<m(s,q)$.

Suppose first that $q<t$. If $t$ is above the line $sq$, then $F,q,t$
is an $n$-cup. If it is below that line, then $q$ is above the chord
$st$, so $P$ and $(s,q,t)$ are opposite chains with common endpoints
and yield an $n$-gon. In the latter case $q$ is not internal to $P$,
whose internal points lie below $st$.

Now let $q>t$. If $q$ is above $ut$, append $q$ to $P$. Otherwise
$q$ is below $ut$. If $t$ is below $sq$, then $(u,t,q)$ is a cap and
$m(t,q)>m(s,q)>m(q,J_3)$; following $(u,t)$ by $J$ without $s$
gives an $(n-1)$-cap.

It remains that $t$ is above $sq$ and $q$ below $ut$. If $u$ is above
$sq$, then $F,u,t$ is an $n$-cup, using
$m(f,s)<m(s,q)<m(s,u)$ and the cup triple $(s,u,t)$. Otherwise
$u$ is below $sq$. If $q$ is above $vu$, replace the last point $t$
of $P$ by $q$; the resulting $(n-1)$-cup and the cap $(s,t,q)$
give an $n$-gon. Finally, if $q$ is below $vu$, then $(v,u,q)$ is a
cap and $m(u,q)>m(s,q)>m(q,J_3)$, giving an $(n-1)$-cap by following
$(v,u)$ with $J$ without $s$.

General position excludes every line equality in this partition of
cases. When $q>t$, the cap suffix starting at $q$ lies after all of
$P$, and the nonterminal points of $F$ lie before $s$. Thus the chains
displayed in the argument have distinct vertices where required.
\end{proof}

\Needspace{14\baselineskip}
The decision map below lists the eight terminal cases of
\cref{lem:threechain}. Use the first test that holds;
$J'=J\setminus\{s\}$ retains its order.
\begin{center}
\small
\begin{tabular}{ll}
\toprule
First applicable test & Forced configuration\\
\midrule
$(f,s,q)$ is a cap & $(f,J)$: $(n-1)$-cap\\
$q<t$ and $t$ above $sq$ & $(F,q,t)$: $n$-cup\\
$q<t$ & $P\cup(s,q,t)$: convex $n$-gon\\
$q$ above $ut$ & $(P,q)$: $n$-cup\\
$t$ below $sq$ & $(u,t,J')$: $(n-1)$-cap\\
$u$ above $sq$ & $(F,u,t)$: $n$-cup\\
$q$ above $vu$ & $(P\setminus\{t\},q)\cup(s,t,q)$: $n$-gon\\
Otherwise & $(v,u,J')$: $(n-1)$-cap\\
\bottomrule
\end{tabular}
\end{center}

\begin{figure}[ht]
\centering
\begin{tikzpicture}[>=Stealth, every node/.style={font=\small},
 box/.style={draw,rounded corners=2pt,inner sep=5pt,align=center}]
 \node[box] (S) at (0,2.2) {$S$};
 \node[box] (U) at (-1.6,1.2) {$U$};
 \node[box] (L) at (1.6,1.2) {$L$};
 \node[box] (G) at (-2.7,0) {$G$\\junction owners};
 \node[box] (R) at (0,0) {$R=U\setminus G$\\$|R|\le D_n$};
 \node[box] (H) at (1.6,-1.3) {$H=L\cup G$\\$|H|\le C_n$};
 \draw[->] (S)--(U); \draw[->] (S)--(L);
 \draw[->] (U)--(G); \draw[->] (U)--(R);
 \draw[->] (L)--(H); \draw[->] (G.south)|-(H.west);
\end{tikzpicture}
\caption{The residual partition. First $S=U\sqcup L$, then
$U=G\sqcup R$. Transferring $G$ to the balanced part gives
$S=H\sqcup R$. Membership in $G$ uses witnesses inside $U$.}
\label{fig:partition}
\end{figure}

\begin{theorem}[Residual transfer]\label{thm:transfer}
For $n\ge6$, $Q_n\le C_n+D_n$ and consequently
$\ES(n)\le C_n+D_n+2$.
\end{theorem}
\begin{proof}
Let $S$ avoid an $(n-1)$-cap and a convex $n$-gon. The case $|S|\le1$
is immediate. At each point $x$, let $p(x)$ be the point minimizing
the slope of the incident unoriented line. The minimizer is unique
by general position. Partition $S$ into
\[
 U=\set{x:p(x)>x},\qquad L=\set{x:p(x)<x}.
\]
The map $p$ takes $U$ into $L$ and $L$ into $U$. For example, if
$x<y=p(x)<z=p(y)$, then $m(y,z)<m(x,y)$, and the weighted-average
identity gives $m(x,z)<m(x,y)$, a contradiction. The reversed argument
handles the other direction. A cap ending in $U$ extends to its chosen
right neighbor in $L$; a cup starting in $L$ extends to its chosen
left neighbor in $U$. Thus $U$ has no $(n-2)$-cap and $L$ has no
$(n-1)$-cup.

Let $G\subseteq U$ be the points ending an $(n-2)$-cup and starting
an $(n-3)$-cap, both in $U$. Put $H=L\cup G$ and $R=U\setminus G$.
The set $R$ satisfies the definition of $D_n$, since any forbidden
junction inside $R$ would also occur in $U$.

Suppose that $H$ contains an $(n-1)$-cup with endpoints $s,t$. If
$s\in L$, extend it left to an $n$-cup. If $t\in G$, extend its
$(n-3)$-cap in $U$ once into $L$, giving an $(n-2)$-cap starting at
$t$. At its junction with the supposed cup, one orientation yields
an $n$-cup and the other an $(n-1)$-cap. Hence $s\in G$ and $t\in L$.
Take the actual $(n-2)$-cup in $U$ ending at $s$ and extend its actual
outgoing $(n-3)$-cap once into $L$. The second point $q$ of this cap
lies in its original cap in $U$; extending that cap only appends its
last point. Therefore $q\ne t$.
\Cref{lem:threechain} applies and gives a contradiction. Hence $H$
has no $(n-1)$-cup, and inherits the other avoidances, so $|H|\le C_n$.
Now $S=H\sqcup R$ gives the result. The threshold statement follows
from \cref{prop:shift}.
\end{proof}
